% =============================================================================
% documentacion_tecnica_final.tex
% Documentación técnica definitiva del proyecto hermetic-bands
% Basada en: Raman & Fan, PRL 104, 087401 (2010)
% Revisión: incluye correcciones de auditoría de código
% =============================================================================

\documentclass[12pt,a4paper]{article}

% ---------------------------------------------------------------------------
% Codificación, idioma y tipografía
% ---------------------------------------------------------------------------
\usepackage[utf8]{inputenc}
\usepackage[T1]{fontenc}
\usepackage[spanish]{babel}
\usepackage{lmodern}
\usepackage{microtype}

% ---------------------------------------------------------------------------
% Geometría de página
% ---------------------------------------------------------------------------
\usepackage[top=2.5cm, bottom=2.5cm, left=2.5cm, right=2.5cm]{geometry}

% ---------------------------------------------------------------------------
% Matemáticas
% ---------------------------------------------------------------------------
\usepackage{amsmath}
\usepackage{amssymb}
\usepackage{amsthm}
\usepackage{bm}

% ---------------------------------------------------------------------------
% Gráficos y diagramas
% ---------------------------------------------------------------------------
\usepackage{graphicx}
\usepackage{tikz}
\usetikzlibrary{shapes.geometric, arrows.meta, positioning, fit, backgrounds}

% ---------------------------------------------------------------------------
% Tablas de publicación científica
% ---------------------------------------------------------------------------
\usepackage{booktabs}
\usepackage{array}
\usepackage{multirow}

% ---------------------------------------------------------------------------
% Bloques de código fuente
% ---------------------------------------------------------------------------
\usepackage{listings}
\usepackage{xcolor}

\definecolor{codebg}{HTML}{F5F5F5}
\definecolor{codeframe}{HTML}{CCCCCC}
\definecolor{keyword}{HTML}{0033B3}
\definecolor{comment}{HTML}{808080}
\definecolor{string}{HTML}{067D17}
\definecolor{number}{HTML}{1750EB}

\lstdefinestyle{cpp}{
  language=C++,
  backgroundcolor=\color{codebg},
  frame=single,
  framesep=4pt,
  rulecolor=\color{codeframe},
  basicstyle=\ttfamily\small,
  keywordstyle=\color{keyword}\bfseries,
  commentstyle=\color{comment}\itshape,
  stringstyle=\color{string},
  numberstyle=\tiny\color{number},
  numbers=left,
  numbersep=8pt,
  showstringspaces=false,
  breaklines=true,
  tabsize=4,
  captionpos=b,
}

\lstdefinestyle{bash}{
  language=bash,
  backgroundcolor=\color{codebg},
  frame=single,
  framesep=4pt,
  rulecolor=\color{codeframe},
  basicstyle=\ttfamily\small,
  keywordstyle=\color{keyword}\bfseries,
  commentstyle=\color{comment}\itshape,
  stringstyle=\color{string},
  showstringspaces=false,
  breaklines=true,
  tabsize=2,
  captionpos=b,
}

\lstdefinestyle{python}{
  language=Python,
  backgroundcolor=\color{codebg},
  frame=single,
  framesep=4pt,
  rulecolor=\color{codeframe},
  basicstyle=\ttfamily\small,
  keywordstyle=\color{keyword}\bfseries,
  commentstyle=\color{comment}\itshape,
  stringstyle=\color{string},
  showstringspaces=false,
  breaklines=true,
  tabsize=4,
  captionpos=b,
}

% ---------------------------------------------------------------------------
% Hiperenlaces
% ---------------------------------------------------------------------------
\usepackage{hyperref}
\hypersetup{
  colorlinks=true,
  linkcolor=blue!60!black,
  citecolor=green!50!black,
  urlcolor=blue!70!black,
}

% ---------------------------------------------------------------------------
% Encabezados y pie de página
% ---------------------------------------------------------------------------
\usepackage{fancyhdr}
\pagestyle{fancy}
\fancyhf{}
\rhead{\small hermetic-bands}
\lhead{\small Documentación Técnica}
\cfoot{\thepage}
\renewcommand{\headrulewidth}{0.4pt}

% ---------------------------------------------------------------------------
% Macros matemáticos
% ---------------------------------------------------------------------------
\newcommand{\R}{\mathbb{R}}
\newcommand{\C}{\mathbb{C}}
\newcommand{\Hhat}{\hat{H}}
\newcommand{\sqA}{\sqrt{A}}
\newcommand{\isqA}{A^{-1/2}}
\newcommand{\omegap}{\omega_p}
\newcommand{\epsinf}{\varepsilon_\infty}
\newcommand{\vect}[1]{\bm{#1}}
\newcommand{\abs}[1]{\left|#1\right|}
\newcommand{\norm}[1]{\left\|#1\right\|}
\newcommand{\conj}[1]{\overline{#1}}
\DeclareMathOperator{\diag}{diag}
\newcommand{\parti}[2]{\frac{\partial #1}{\partial #2}}

% ---------------------------------------------------------------------------
% Ajustes de espaciado
% ---------------------------------------------------------------------------
\setlength{\parskip}{0.4em}
\setlength{\parindent}{1.5em}

% =============================================================================
\begin{document}
% =============================================================================

% ---------------------------------------------------------------------------
% Portada
% ---------------------------------------------------------------------------
\begin{titlepage}
  \centering
  \vspace*{2cm}
  {\Large\bfseries Universidad de Medellín\\[0.3cm]
  Facultad de Ciencias Básicas\\[0.3cm]
  Física Computacional --- Período 6\par}
  \vspace{2cm}
  \rule{\linewidth}{0.5pt}\\[0.4cm]
  {\LARGE\bfseries Estructura de Bandas Fotónicas de\\[0.3cm]
  Metamateriales Dispersivos:\\[0.3cm]
  Formulación Hermitiana con Variable Auxiliar\par}
  {\large Documentación Técnica del Proyecto \texttt{hermetic-bands}\\
  Versión definitiva con correcciones de auditoría\par}
  \rule{\linewidth}{0.5pt}\\[1.5cm]
  {\large\bfseries Autor:} Emanuel Cabrera Novoa\\[0.3cm]
  {\large\bfseries Correo:} lacabrazapa@gmail.com\\[0.3cm]
  {\large\bfseries Fecha:} Mayo 2026\\[2cm]
  \begin{abstract}
    \noindent
    Este documento describe la implementación computacional del método de
    Raman y Fan~\cite{Raman2010} para el cálculo de estructuras de bandas
    fotónicas en cristales plasmónicos dispersivos.
    La formulación transforma el problema de autovalores no-Hermitiano
    $\omega^2 \varepsilon(\omega) \vect{E} = c^{-2}\nabla\times\nabla\times\vect{E}$
    en el problema Hermitiano definido-positivo
    $\omega\, \hat{H}\, \hat{y} = \omega\, \hat{y}$
    mediante la introducción de un campo auxiliar vectorial $\vect{V}$
    que absorbe la dinámica del modelo de Drude.
    Se detallan la discretización mediante la malla de Yee estacionada,
    la construcción de operadores de rotacional con condiciones de contorno
    de Bloch, el ensamblado en formato CSR, y la resolución mediante
    ARPACK con inversión espectral (shift-invert) y factorización LU densa
    (LAPACK \texttt{zgetrf}/\texttt{zgetrs}).
    Se documenta además la corrección de dos errores críticos identificados
    en la auditoría del código: la doble transformación espectral en
    \texttt{zneupd} (modo 3) y la selección errónea del criterio
    \texttt{"LM"} de ARPACK.
  \end{abstract}
\end{titlepage}

\tableofcontents
\newpage

% =============================================================================
\section{Introducción}
% =============================================================================

Los metamateriales plasmónicos son estructuras periódicas que contienen
inclusiones metálicas nanométricas con permitividad dieléctrica descrita
por el modelo de Drude:
\begin{equation}
  \varepsilon(\omega) = \epsinf\!\left(1 - \frac{\omegap^2}{\omega^2 + i\Gamma\omega}\right),
  \label{eq:drude}
\end{equation}
donde $\epsinf$ es la permitividad a alta frecuencia, $\omegap$ la
frecuencia de plasma y $\Gamma$ la tasa de amortiguamiento.
La dependencia en frecuencia de~\eqref{eq:drude} hace que las ecuaciones
de Maxwell resulten en un problema de autovalores polinomial en $\omega$,
que destruye la hermiticidad del operador Maxwell estándar y vuelve
ineficientes los métodos clásicos de cálculo de bandas fotónicas.

Raman y Fan~\cite{Raman2010} resuelven esta dificultad introduciendo un
campo auxiliar $\vect{V}(\vect{r},t)$ que satisface la ecuación de movimiento
de los electrones libres en el metal.
La idea central es reformular el problema cuadrático en $\omega$ como un
problema lineal de autovalores $\omega\, A\, x = B\, x$, donde tanto $A$
como $B$ son matrices Hermitian as definidas-positivas.
La transformación de similitud $\hat{H} = A^{-1/2} B A^{-1/2}$ convierte
este sistema en la ecuación canónica $\omega\,\hat{y} = \hat{H}\,\hat{y}$,
resoluble eficientemente con el método de Lanczos--Arnoldi
implícitamente reiniciado (IRLM) implementado en ARPACK.

El proyecto \texttt{hermetic-bands} implementa esta formulación en C++17
sobre una malla de Yee 2D con condiciones de contorno de Bloch,
soportando los modos de polarización TM y TE, inclusiones cuadradas y
circulares, el cálculo de la estructura de bandas a lo largo de la ruta
$\Gamma\!\to\! X\!\to\! M\!\to\!\Gamma$ de la primera zona de Brillouin,
y correcciones perturbativas de primer orden para los efectos de pérdidas.

% =============================================================================
\section{Marco Teórico}
% =============================================================================

\subsection{Ecuaciones de Maxwell y convención temporal}

Adoptamos la convención $e^{-i\omega t}$ para todos los campos
armónico-temporales:
\begin{align}
  \nabla \times \vect{E} &= +i\omega\mu_0\vect{H}
    \qquad \text{(Faraday)}, \label{eq:faraday}\\
  \nabla \times \vect{H} &= -i\omega\varepsilon(\omega)\vect{E}
    \qquad \text{(Ampère--Maxwell)}. \label{eq:ampere}
\end{align}
Combinando~\eqref{eq:faraday} y~\eqref{eq:ampere} con el modelo de
Drude~\eqref{eq:drude} se obtiene el problema de autovalores implícito en
$\omega^2$, que es no-Hermitiano para $\Gamma>0$ e incluso para
$\Gamma=0$ cuando se emplea el modelo de Drude sin pérdidas.

\subsection{Variable auxiliar de Drude}

Definimos el campo auxiliar $\vect{V}(\vect{r},t) = -i\omega\vect{P}/(\epsinf\omegap^2)$,
donde $\vect{P}$ es la polarización del plasma.
La ecuación de movimiento de los electrones libres con amortiguamiento $\Gamma$ es:
\begin{equation}
  (-i\omega + \Gamma)\,\vect{V} = \epsinf\omegap^2\,\vect{E}.
  \label{eq:aux_field}
\end{equation}
Para el caso sin pérdidas ($\Gamma=0$), la ecuación~\eqref{eq:aux_field}
se convierte en $-i\omega\vect{V} = \epsinf\omegap^2\vect{E}$,
que combinada con Ampère permite eliminar la dependencia implícita en $\omega$.

\subsection{Reformulación de Raman y Fan (ec.\ 6)}

\label{sec:raman_fan}

El vector de estado expandido para el modo \textbf{TM} es
$x = (\vect{H}_x,\, \vect{H}_y,\, \vect{E}_z,\, \vect{V}_z)^T$,
donde cada bloque tiene $N^2$ componentes sobre la malla $N\times N$.
El sistema lineal de autovalores resultante es~\cite{Raman2010}:
\begin{equation}
  \omega\, A\, x = B\, x,
  \label{eq:AxBx}
\end{equation}
con la matriz de masa $A = \diag(A_{Hx}, A_{Hy}, A_{Ez}, A_{Vz})$:
\begin{align}
  A_{H\alpha}[k] &= \mu_0 \quad \forall k, \\
  A_{Ez}[k]      &= \epsinf(r_k), \\
  A_{Vz}[k]      &=
    \begin{cases}
      \dfrac{1}{\epsinf(r_k)\,\omegap^2(r_k)} & k \in \text{metal},\\[6pt]
      1 & k \in \text{aire (regularización)},
    \end{cases}
  \label{eq:A_diag}
\end{align}
y la matriz dinámica $B$ en bloques:
\begin{equation}
  B =
  \begin{pmatrix}
    0      &  0      & -i C_E^{TM} & 0  \\
    0      &  0      &     0       & 0  \\
    +i C_H^{TM} & 0  &  0         & -i I_m \\
    0      &  0      & +i I_m      &  0
  \end{pmatrix},
  \label{eq:B_TM}
\end{equation}
donde $C_E^{TM}: \R^{N^2}\to\R^{2N^2}$ es el operador de rotacional
discreto de $\vect{E}$ e $I_m$ denota la identidad restringida a los
nodos metálicos.
La estructura de bloques para el modo \textbf{TE} con vector
$x=(\vect{H}_z, \vect{E}_x, \vect{E}_y, \vect{V}_x, \vect{V}_y)^T$
es análoga con cinco bloques.

\subsection{Transformación de similitud y Hamiltoniano Hermitiano}

Dado que $A$ es diagonal definida-positiva, su raíz cuadrada
$S = A^{1/2} = \diag(\sqrt{A_{kk}})$ se calcula en $\mathcal{O}(N^2)$
operaciones.
El cambio de variable $\hat{y} = S\, x$ transforma~\eqref{eq:AxBx} en:
\begin{equation}
  \omega\,\hat{y} = \Hhat\,\hat{y},
  \qquad
  \Hhat = S^{-1} B S^{-1} = \isqA B \isqA.
  \label{eq:Hhat}
\end{equation}
La hermiticidad de $\Hhat$ se hereda de $B$ siempre que $C_H = C_E^\dagger$,
condición que se verifica en la siguiente sección.
Computacionalmente, la transformación~\eqref{eq:Hhat} se aplica
directamente sobre los elementos no-cero del formato CSR:
\begin{equation}
  \Hhat_{ij} = s^{-1}_i\, B_{ij}\, s^{-1}_j,
  \label{eq:similarity_sparse}
\end{equation}
lo que preserva el patrón de no-ceros de $B$ y requiere $\mathcal{O}(\text{nnz})$
operaciones, sin necesidad de construir $A^{-1/2}$ como matriz completa.

\subsection{Discretización de Yee y hermiticidad de los operadores de rotacional}

\label{sec:yee}

La celda unitaria $[0,a)\times[0,a)$ se discretiza en $N\times N$ celdas
de lado $\Delta = a/N$.
Las posiciones de los campos en el modo TM son:
\begin{align*}
  E_z[i,j] &: \text{ nodo primario } (i\Delta,\, j\Delta),\\
  H_x[i,j] &: \text{ arista-}y \;\; (i\Delta,\, (j+\tfrac{1}{2})\Delta),\\
  H_y[i,j] &: \text{ arista-}x \;\; ((i+\tfrac{1}{2})\Delta,\, j\Delta).
\end{align*}

Los operadores de diferencias finitas hacia adelante y hacia atrás son:
\begin{equation}
  (D_f^x\, f)[i,j] = \frac{f[i+1,j] - f[i,j]}{\Delta},
  \qquad
  (D_b^x\, f)[i,j] = \frac{f[i,j] - f[i-1,j]}{\Delta},
  \label{eq:Dfx}
\end{equation}
con condiciones periódicas de Bloch al cruzar la frontera de la celda:
\begin{equation}
  f[N,j] = e^{+ik_x a}\, f[0,j], \qquad
  f[-1,j] = e^{-ik_x a}\, f[N-1,j].
  \label{eq:bloch}
\end{equation}

\textbf{Hermiticidad del par $(C_E, C_H)$.}
El rotacional discreto de Faraday (diferencias hacia adelante) y el de
Ampère (diferencias hacia atrás) satisfacen:
\begin{equation}
  C_H^{TM} = \bigl(C_E^{TM}\bigr)^\dagger.
  \label{eq:curl_adjoint}
\end{equation}
La verificación de~\eqref{eq:curl_adjoint} \emph{no es trivial} en
presencia de las fases de Bloch.
Para el operador $D_f^x$, la entrada fuera de la diagonal asociada al
cruce de frontera tiene valor $+e^{+ik_x}/\Delta$;
su adjunto conjugado, que corresponde a $D_b^x$, tiene el valor
$-e^{-ik_x}/\Delta$ en la posición transpuesta.
Es decir, los factores de Bloch \emph{deben cambiar de signo en el
exponente al transponerlos}:
\begin{equation}
  \bigl[(D_f^x)^\dagger\bigr]_{(i-1),i}
    = \conj{(D_f^x)_{i,(i-1)}}
    = \conj{\!\left(\frac{e^{+ik_x}}{\Delta}\right)\!}
    = \frac{e^{-ik_x}}{\Delta},
  \label{eq:bloch_conj}
\end{equation}
que coincide exactamente con la entrada de $D_b^x$.
En el código, esta propiedad se garantiza por construcción:
\begin{lstlisting}[style=cpp, caption={Construcción del adjunto exacto en \texttt{operators.cpp}}]
// C_H_TM = C_E_TM dagger -- garantiza la relacion de adjunto exacta.
CSRMatrix build_CH_TM(const YeeGrid& g, Real kx, Real ky) {
    return conj_transpose(build_CE_TM(g, kx, ky));
}
\end{lstlisting}

Cualquier asimetría residual $\norm{\Hhat - \Hhat^\dagger}_F$ introduce
autovalores espurios con parte imaginaria no nula que contaminarían el
espectro físico.
La función \texttt{check\_hermitian()} de \texttt{sparse.hpp} calcula el
máximo elemento de esta diferencia y puede invocarse en modo de diagnóstico
(\texttt{-{}-verbosity 2}).

% =============================================================================
\section{Implementación Numérica}
% =============================================================================

\subsection{Almacenamiento disperso CSR}

Todos los operadores ($D_f^x$, $D_b^x$, $C_E$, $C_H$, $B$, $\Hhat$) se
almacenan en formato \emph{Compressed Sparse Row} (CSR):
\begin{itemize}
  \item \texttt{val[k]}: valor del no-cero $k$.
  \item \texttt{col\_ind[k]}: índice de columna del no-cero $k$.
  \item \texttt{row\_ptr[i]}: índice en \texttt{val} del primer no-cero de la fila $i$.
\end{itemize}
El producto matriz-vector $y \leftarrow \alpha A x + \beta y$ recorre
\texttt{val[]} y \texttt{col\_ind[]} secuencialmente, lo que es
favorable para la caché.
La densidad de no-ceros por fila es $\leq 5$ para los operadores de
diferencias finitas en 2D, y $\leq 10$ para la matriz dinámica $B$.

\subsection{Inversión espectral y factorización LU densa}
\label{sec:shift_invert}

ARPACK en modo estándar (\texttt{mode}=1) encuentra los autovalores de
\emph{mayor magnitud}, que para un Hamiltoniano fotónico corresponden a
los modos de alta frecuencia.
La \emph{inversión espectral} (shift-invert) reemplaza el problema
$\omega\,\hat{y} = \Hhat\,\hat{y}$ por:
\begin{equation}
  \lambda\,\hat{y} = (\Hhat - \sigma I)^{-1}\hat{y},
  \qquad
  \lambda = \frac{1}{\omega - \sigma},
  \label{eq:shift_invert}
\end{equation}
de modo que solicitar los $\lambda$ de mayor módulo recupera los $\omega$
más cercanos al desplazamiento $\sigma$.

\textbf{Selección del desplazamiento.}
Eligiendo $0 < \sigma < \omega_{\min}^{\text{físico}}$, los modos del
espacio nulo ($\omega=0$, provenientes de los grados de libertad $\vect{V}$
en el aire) mapean a $\lambda = -1/\sigma < 0$, mientras que los modos
físicos con $\omega > \sigma$ mapean a $\lambda > 0$.
Por defecto se usa $\sigma = 0{,}01\,\omegap$; el usuario puede sobreescribir
con el flag \texttt{-{}-sigma}.

\textbf{Justificación del LU denso.}
En cada paso de Lanczos, ARPACK solicita la aplicación de
$(\Hhat - \sigma I)^{-1}$.
Para el rango de validación del artículo original ($N \leq 30$,
$n_\text{dof} \leq 3600$ para TM), se emplea la factorización LU densa
de LAPACK (\texttt{zgetrf}/\texttt{zgetrs}):
\begin{align}
  \text{Factorización (una vez):}\quad &
  \mathcal{O}(n_\text{dof}^3) \approx 10^{10}\;\text{flops para } N=30. \\
  \text{Resolución (por iteración):}\quad &
  \mathcal{O}(n_\text{dof}^2) \approx 1{,}3\times10^7\;\text{flops para } N=30.
\end{align}
Esta elección proporciona:
\begin{enumerate}
  \item \textbf{Determinismo absoluto}: el resultado es bit a bit
    reproducible entre plataformas, esencial para verificar figuras del
    artículo.
  \item \textbf{Precisión de máquina}: el error de retroceso del LU es
    $\lesssim 10^{-13}$, muy por debajo del error de discretización
    ($\sim 1\%$ a $N=20$).
  \item \textbf{Sin dependencias adicionales}: no se requiere UMFPACK,
    PARDISO ni SuperLU; solo BLAS + LAPACK ya declarados como dependencias.
\end{enumerate}

\textbf{Límite de memoria.}
El costo de memoria de la factorización escala como $n_\text{dof}^2 \times 16\,\text{B}$:
\begin{equation}
  M_\text{LU} \approx (4N^2)^2 \times 16\,\text{B} = 256\,N^4\,\text{B} \quad \text{(TM)}.
\end{equation}
Para $N=40$ esto asciende a $\approx 655\,\text{MB}$; para $N=50$,
$\approx 1{,}6\,\text{GB}$.
El código emite una advertencia automática cuando $N > 40$, indicando que
para resoluciones mayores conviene reemplazar el solucionador interno por
un solver disperso directo (PARDISO, MUMPS) o un método iterativo
precondicionado (GMRES + ILU).

\subsection{Bucle de comunicación inversa de ARPACK}

El bucle de comunicación inversa de \texttt{znaupd} tiene la siguiente
estructura:

\begin{lstlisting}[style=cpp, caption={Bucle principal ARPACK en \texttt{eigensolver.cpp}}]
while (true) {
    znaupd_(&ido, bmat, &N, which, &nev, &tol,
            resid.data(), &ncv, v.data(), &N,
            iparam, ipntr,
            workd.data(), workl.data(), &lworkl,
            rwork.data(), &info);

    if (ido == -1 || ido == 1) {
        // Aplicar y = (H_hat - sigma*I)^{-1} * x
        op.apply(x_in, y_out);   // usa zgetrs internamente
    }
    else if (ido == 2) {
        // B = I: copiar x -> y directamente
        std::copy(workd.begin() + x_off,
                  workd.begin() + x_off + N,
                  workd.begin() + y_off);
    }
    else if (ido == 99) break;   // convergido
}
\end{lstlisting}

\subsection{Correcciones de bugs críticos identificados en la auditoría}

\label{sec:bugs}

\subsubsection{Bug 1 --- Doble transformación espectral en \texttt{zneupd}}

Tras la convergencia del bucle de \texttt{znaupd}, la subrutina
\texttt{zneupd} extrae los pares (autovalor, autovector) del espacio
de Krylov.
En \textbf{modo 3 (SHIFTI)}, \texttt{zneupd} aplica internamente la
transformación inversa:
\begin{equation}
  d_k \leftarrow \sigma + \frac{1}{\rho_k},
  \qquad \rho_k = \text{valor de Ritz en el espacio transformado},
  \label{eq:zneupd_backtransform}
\end{equation}
de modo que el arreglo \texttt{d[]} a la salida de \texttt{zneupd}
\emph{ya contiene los autovalores $\omega$ del problema original}.

La versión defectuosa del código aplicaba nuevamente la conversión
$\sigma + 1/\lambda$, produciendo valores completamente erróneos:
\begin{lstlisting}[style=cpp, caption={Código defectuoso (ELIMINADO)}]
// ERROR: zneupd ya aplico sigma + 1/rho internamente.
// Esta segunda conversion da valores incorrectos.
const Complex lambda = d[j];
const Complex omega  = Complex{sigma, 0.0} + 1.0 / lambda;
result.eigenvalues.push_back(omega.real());
\end{lstlisting}

La corrección lee directamente la parte real de \texttt{d[j]}:
\begin{lstlisting}[style=cpp, caption={Código corregido en \texttt{eigensolver.cpp}}]
// zneupd (modo 3, SHIFTI) aplica d[k] = sigma + 1/rho[k] internamente.
// d[] ya contiene los autovalores omega del problema original.
for (int j = 0; j < nconv; ++j) {
    result.eigenvalues.push_back(
        d[static_cast<std::size_t>(j)].real());
}
\end{lstlisting}

\subsubsection{Bug 2 --- Criterio de selección \texttt{"LM"} en ARPACK}

El parámetro \texttt{which} de \texttt{znaupd} indica qué autovalores
del operador $(\Hhat - \sigma I)^{-1}$ se solicitan.
Con el desplazamiento shift-invert, los modos físicos ($\omega > \sigma$)
mapean a valores de Ritz $\lambda = 1/(\omega-\sigma) > 0$, mientras
que los modos del espacio nulo ($\omega = 0$) mapean a
$\lambda = -1/\sigma < 0$.

La selección \texttt{"LM"} (\emph{Largest Magnitude}) toma los $|\lambda|$
más grandes, y dado que $|-1/\sigma| = 1/\sigma \gg 1$ para $\sigma$
pequeño, ARPACK podía seleccionar preferentemente los modos no físicos
del espacio nulo.
La corrección usa \texttt{"LR"} (\emph{Largest Real part}), que excluye
limpiamente todos los $\lambda < 0$:

\begin{lstlisting}[style=cpp, caption={Corrección del criterio ARPACK en \texttt{eigensolver.cpp}}]
// "LR" (mayor parte real) selecciona modos fisicos (lambda > 0).
// "LM" (mayor magnitud) podia incluir modos del espacio nulo
//  con lambda = -1/sigma << 0 de gran magnitud absoluta.
char which[] = "LR";
\end{lstlisting}

\subsection{Tratamiento de pérdidas}

\label{sec:losses}

Con amortiguamiento $\Gamma \neq 0$, la ecuación de movimiento~\eqref{eq:aux_field}
adquiere un término extra en la diagonal del bloque $V$:
\begin{equation}
  \delta B_{VV}[k] = -i\Gamma\cdot\mathbf{1}_{k\in\text{metal}},
  \qquad
  \delta\Hhat_{VV}[k] = \frac{-i\Gamma}{\epsinf(r_k)\,\omegap^2(r_k)}.
  \label{eq:delta_H}
\end{equation}
Esta perturbación es anti-Hermitiana ($\delta\Hhat^\dagger = -\delta\Hhat$),
de modo que los autovalores del sistema perturbado son complejos:
$\omega = \omega_r + i\omega_i$ con $\omega_i < 0$ (decaimiento físico).

\textbf{Separación lossless/lossy en el código.}
La perturbación~\eqref{eq:delta_H} entra únicamente en
\texttt{assemble\_nonhermitian()}, que construye $\Hhat_\text{nh} = \Hhat_0 + \delta\Hhat$.
El Hamiltoniano sin pérdidas $\Hhat_0$ permanece intacto, garantizando que
los autovalores de referencia sean reales.

\textbf{Corrección perturbativa (ec.\ 15 de Raman y Fan).}
Para $\Gamma \ll \omegap$, la corrección de primer orden a la frecuencia
del modo $n$ es:
\begin{equation}
  \omega_n^{(1)} = \langle \hat{y}_n^{(0)} | \delta\Hhat | \hat{y}_n^{(0)} \rangle
    = \frac{-i\Gamma\displaystyle\int_\text{metal}
        \frac{|V_n|^2}{\epsinf\,\omegap^2}\,d^2r}
      {\displaystyle\int \Bigl(\mu_0|H_n|^2 + \epsinf|E_n|^2
        + \frac{|V_n|^2}{\epsinf\omegap^2}\Bigr)d^2r}.
  \label{eq:perturbative}
\end{equation}
Esta expresión usa los modos $\hat{y}_n^{(0)}$ del sistema sin pérdidas,
evitando re-resolver el problema no-Hermitiano completo.

% =============================================================================
\section{Arquitectura del Software}
% =============================================================================

\subsection{Diagrama de flujo de datos}

La figura~\ref{fig:flowchart} muestra el flujo de datos principal.

\begin{figure}[htbp]
  \centering
  \begin{tikzpicture}[
    node distance=1.3cm and 2.0cm,
    block/.style={
      rectangle, rounded corners=4pt,
      draw=blue!60!black, fill=blue!8,
      text width=4.2cm, align=center,
      minimum height=0.9cm, font=\small
    },
    decision/.style={
      diamond, draw=orange!70!black, fill=orange!12,
      text width=2.8cm, align=center,
      inner sep=1pt, font=\small
    },
    io/.style={
      trapezium, trapezium left angle=70, trapezium right angle=110,
      draw=green!60!black, fill=green!10,
      text width=3.6cm, align=center,
      minimum height=0.8cm, font=\small
    },
    arr/.style={-Stealth, thick, blue!50!black},
  ]

  \node[io]     (cli)   {Interfaz CLI\\(\texttt{parse\_args})};
  \node[block]  (param) [below=of cli]
                        {Parámetros\\(\texttt{SimParams})};
  \node[block]  (grid)  [below=of param]
                        {Malla de Yee\\(\texttt{YeeGrid})};
  \node[block]  (assem) [below=of grid]
                        {Ensamblado $\hat{H}$ CSR\\(\texttt{assemble\_hermitian})};
  \node[block]  (lu)    [below=of assem]
                        {Factor. LU densa\\(\texttt{ShiftInvertOp::factor})};
  \node[block]  (arpack)[below=of lu]
                        {ARPACK \texttt{znaupd/zneupd}\\(shift-invert, \texttt{"LR"})};
  \node[block]  (back)  [below=of arpack]
                        {Back-transform.\\$x = S^{-1}\hat{y}$};
  \node[io]     (out)   [below=of back]
                        {Salida \texttt{.dat}\\(\texttt{write\_bands})};

  \node[decision](loss) [right=2.2cm of arpack]
                        {$\Gamma > 0$?};
  \node[block]  (nh)    [right=2.2cm of loss]
                        {Sistema NH\\+ pérdidas exactas};
  \node[block]  (pert)  [below=1.1cm of nh]
                        {Corrección\\perturbativa (ec.\ 15)};

  \draw[arr] (cli)    -- (param);
  \draw[arr] (param)  -- (grid);
  \draw[arr] (grid)   -- (assem);
  \draw[arr] (assem)  -- (lu);
  \draw[arr] (lu)     -- (arpack);
  \draw[arr] (arpack) -- (back);
  \draw[arr] (back)   -- (out);
  \draw[arr] (arpack) -- (loss);
  \draw[arr] (loss)   -- node[above,font=\tiny]{Sí} (nh);
  \draw[arr] (nh)     -- (pert);
  \draw[arr] (loss.south) -- node[right,font=\tiny]{No}
             ++(0,-0.7) -| (back.east);

  \end{tikzpicture}
  \caption{Flujo de datos del proyecto \texttt{hermetic-bands}.
    El camino principal (izquierda) ejecuta el cálculo sin pérdidas.
    La ramificación derecha se activa cuando $\Gamma > 0$.}
  \label{fig:flowchart}
\end{figure}

\subsection{Módulos del proyecto}

\begin{table}[htbp]
  \centering
  \caption{Módulos C++ del proyecto y su responsabilidad.}
  \label{tab:modules}
  \begin{tabular}{@{}lll@{}}
    \toprule
    \textbf{Archivo} & \textbf{Módulo} & \textbf{Responsabilidad} \\
    \midrule
    \texttt{utils.hpp/.cpp}       & Tipos y CLI & Tipos base, constantes, parseo de argumentos \\
    \texttt{sparse.hpp/.cpp}      & Motor CSR   & Almacenamiento, matvec, adjunto conjugado \\
    \texttt{yee\_grid.hpp/.cpp}   & Malla Yee   & Geometría, máscaras de material, promediado armónico \\
    \texttt{operators.hpp/.cpp}   & Operadores  & $C_E$, $C_H$, $A$, $B$, $\hat{H}$ \\
    \texttt{matrices.hpp/.cpp}    & Sistemas    & Ensamblado de \texttt{HermitianSystem} y \texttt{NonHermitianSystem} \\
    \texttt{eigensolver.hpp/.cpp} & Solver      & Bucle ARPACK, back-transformación, residuos \\
    \texttt{perturbation.hpp/.cpp}& Pérdidas    & Ec.~(15) perturbativa y ortonormalidad \\
    \texttt{output.hpp/.cpp}      & Salida      & Escritura de archivos \texttt{.dat} \\
    \texttt{main.cpp}             & Orquestador & Flujo principal de la simulación \\
    \bottomrule
  \end{tabular}
\end{table}

% =============================================================================
\section{Referencia Completa de la Interfaz de Línea de Comandos}
% =============================================================================

\subsection{Compilación}

\begin{lstlisting}[style=bash, caption={Compilación en modo Release (recomendado)}]
# Desde el directorio raiz del proyecto
mkdir -p build_cmake
cmake -DCMAKE_BUILD_TYPE=Release -B build_cmake .
cmake --build build_cmake -j$(nproc)
# Binario resultante: build_cmake/hermetic-bands
\end{lstlisting}

Las banderas de optimización activadas en modo Release son
\texttt{-O3 -march=native -ffast-math} para C++ y
\texttt{-O3} para los módulos Fortran de ARPACK.

\subsection{Tabla de parámetros}

\begin{table}[htbp]
  \centering
  \caption{Parámetros de la interfaz de línea de comandos (CLI).
    Los valores entre corchetes son los valores por defecto.}
  \label{tab:cli}
  \small
  \begin{tabular}{@{}lll@{}}
    \toprule
    \textbf{Flag} & \textbf{Tipo / Rango} & \textbf{Descripción} \\
    \midrule
    \texttt{-{}-mode [TE|TM]}          & enumeración     & Polarización del modo [TE] \\
    \texttt{-{}-resolution N}          & entero $\geq 4$ & Puntos Yee por lado de la celda [20] \\
    \texttt{-{}-fill-fraction f}       & real $(0,1)$    & Fracción de llenado metálico [0.25] \\
    \texttt{-{}-shape [square|circle]} & enumeración     & Geometría de la inclusión [square] \\
    \texttt{-{}-eps-inf VALUE}         & real $>0$       & Permitividad $\epsinf$ [1.0] \\
    \texttt{-{}-nk N}                  & entero $\geq 2$ & Puntos $k$ en la ruta [20] \\
    \texttt{-{}-nbands N}              & entero $\geq 1$ & Número de bandas a calcular [15] \\
    \texttt{-{}-sigma S}               & real $\geq 0$   & Desplazamiento ARPACK; 0=auto [0] \\
    \texttt{-{}-gamma G}               & real $\geq 0$   & Tasa de amortiguamiento $\Gamma/\omegap$ [0] \\
    \texttt{-{}-perturb}               & booleano        & Calcular corrección perturbativa \\
    \texttt{-{}-output FILE}           & cadena          & Archivo de salida [bands.dat] \\
    \texttt{-{}-field-mode N}          & entero $\geq 0$ & Índice de banda para exportar campo \\
    \texttt{-{}-field-kindex I}        & entero $\geq 0$ & Índice de punto-$k$ para el campo [0] \\
    \texttt{-{}-verify-orthogonality}  & booleano        & Verificar ortonormalidad (ec.\ 13) \\
    \texttt{-{}-convergence}           & booleano        & Ejecutar estudio de convergencia \\
    \texttt{-{}-conv-N-start N}        & entero          & Resolución inicial de convergencia [8] \\
    \texttt{-{}-conv-N-end N}          & entero          & Resolución final de convergencia [30] \\
    \texttt{-{}-conv-N-step N}         & entero          & Paso de resolución [4] \\
    \texttt{-{}-full-bz}               & booleano        & Ruta $\Gamma\to X\to M\to\Gamma$ completa \\
    \texttt{-{}-map2d}                 & booleano        & Mapa 2D $\omega(k_x,k_y)$ \\
    \texttt{-{}-map2d-nk N}            & entero $\geq 2$ & Puntos por eje en el mapa 2D [10] \\
    \texttt{-{}-verbosity [0|1|2]}     & entero          & Nivel de salida en consola [1] \\
    \bottomrule
  \end{tabular}
\end{table}

\noindent
\textbf{Notas sobre flags deprecados (compatibilidad hacia atrás):}
\begin{itemize}
  \item \texttt{-{}-loss} es sinónimo de \texttt{-{}-gamma} (deprecado).
  \item \texttt{-{}-verify-ortho} es sinónimo de \texttt{-{}-verify-orthogonality} (deprecado).
\end{itemize}

% =============================================================================
\section{Ejemplos de Uso}
% =============================================================================

\subsection{Ejemplo 1: Estructura de bandas TM estándar}

\begin{lstlisting}[style=bash, caption={Estructura de bandas TM con resolución $30\times30$, 10 bandas.}]
./build_cmake/hermetic-bands \
    --mode TM \
    --resolution 30 \
    --fill-fraction 0.25 \
    --shape square \
    --nk 40 \
    --nbands 10 \
    --output bands_TM.dat

python3 visualize_results.py bands_TM.dat
\end{lstlisting}

El archivo \texttt{bands\_TM.dat} contiene columnas:
\texttt{k\_index}, \texttt{k*a/pi}, $\omega_1/\omegap$, $\omega_2/\omegap$, \ldots

\subsection{Ejemplo 2: Exportación del perfil de campo espacial}

Los flags \texttt{-{}-field-mode} y \texttt{-{}-field-kindex} son
\emph{independientes} a partir de la versión actual (a diferencia de
la sintaxis anterior que aceptaba dos argumentos posicionales consecutivos).

\begin{lstlisting}[style=bash, caption={Exportar el perfil de campo de la banda 3 en el punto $k$ de índice 20.}]
./build_cmake/hermetic-bands \
    --mode TM \
    --resolution 30 \
    --nk 40 \
    --field-mode 3 \
    --field-kindex 20 \
    --output /dev/null

python3 visualize_results.py field_3_k20.dat
\end{lstlisting}

\subsection{Ejemplo 3: Estudio de convergencia numérica}

El flag \texttt{-{}-convergence} es \emph{booleano puro}: no acepta
el argumento \texttt{true} ni ningún valor.

\begin{lstlisting}[style=bash, caption={Estudio de convergencia en el punto $X$ ($k_x=\pi$) de $N=8$ a $N=32$.}]
./build_cmake/hermetic-bands \
    --mode TM \
    --convergence \
    --conv-N-start 8 \
    --conv-N-end 32 \
    --conv-N-step 4 \
    --nbands 5 \
    --output /dev/null

python3 visualize_results.py convergence.dat
\end{lstlisting}

El archivo \texttt{convergence.dat} incluye columnas
$N$, $\omega_1(N)$, $\omega_2(N)$, \ldots y el tiempo de cómputo.
El error relativo debería escalar aproximadamente como $\mathcal{O}(N^{-2})$,
consistente con la exactitud de segundo orden del operador de Yee.

\subsection{Ejemplo 4: Zona de Brillouin completa $\Gamma\to X\to M\to\Gamma$}

\begin{lstlisting}[style=bash, caption={Estructura de bandas en la zona de Brillouin completa.}]
./build_cmake/hermetic-bands \
    --mode TM \
    --full-bz \
    --resolution 20 \
    --nk 60 \
    --nbands 10 \
    --output bands_fullBZ.dat

python3 visualize_results.py bands_fullBZ.dat
\end{lstlisting}

La ruta $\Gamma\to X\to M\to\Gamma$ tiene longitudes de arco
$|\Gamma X| = |\text{XM}| = \pi$ y $|M\Gamma| = \pi\sqrt{2}$.
Los puntos se distribuyen con espaciado uniforme en longitud de arco,
etiquetando el eje horizontal en unidades de $\pi/a$.

\subsection{Ejemplo 5: Inclusiones circulares con pérdidas}

\begin{lstlisting}[style=bash, caption={Estructura de bandas con inclusiones circulares y corrección perturbativa.}]
./build_cmake/hermetic-bands \
    --mode TM \
    --shape circle \
    --fill-fraction 0.30 \
    --resolution 24 \
    --full-bz \
    --nk 40 \
    --nbands 10 \
    --gamma 0.01 \
    --perturb \
    --output bands_circle.dat

python3 visualize_results.py loss_comparison.dat
\end{lstlisting}

Para \texttt{-{}-shape circle}, el parámetro \texttt{fill-fraction}
es el radio normalizado $r/a$.
El archivo \texttt{loss\_comparison.dat} contiene, para cada modo,
la frecuencia lossless $\omega_0$, la frecuencia no-Hermitiana exacta
$\text{Re}[\omega]$ e $\text{Im}[\omega]$, y la corrección perturbativa $\omega_0 + \omega^{(1)}$.

% =============================================================================
\section{Visualización con Python}
% =============================================================================

El script \texttt{visualize\_results.py} detecta el tipo de datos por
el nombre del archivo y despacha a la función apropiada.
El fragmento central para la estructura de bandas es:

\begin{lstlisting}[style=python, caption={Función de graficación de bandas en \texttt{visualize\_results.py}}]
import numpy as np
import matplotlib.pyplot as plt
import os

def plot_bands(filename):
    if not os.path.exists(filename):
        print(f"Archivo {filename} no encontrado.")
        return

    data = np.loadtxt(filename)   # shape: (n_kpoints, 2 + nbands)
    if data.ndim < 2:
        print(f"Datos invalidos en {filename}")
        return

    k     = data[:, 1]     # columna 1: k*a/pi
    bands = data[:, 2:]    # columnas 2..: omega_n / omega_p

    plt.figure(figsize=(8, 6))
    for i in range(bands.shape[1]):
        plt.plot(k, bands[:, i])

    plt.xlabel(r'$k a / \pi$')
    plt.ylabel(r'$\omega / \omega_p$')
    plt.title('Estructura de Bandas Fotonica')
    plt.grid(True, linestyle='--', alpha=0.6)
    out_img = filename.replace('.dat', '.png')
    plt.savefig(out_img, dpi=300, bbox_inches='tight')
    print(f"Imagen guardada: {out_img}")
\end{lstlisting}

\noindent
Para el estudio de convergencia:

\begin{lstlisting}[style=python, caption={Función de convergencia en \texttt{visualize\_results.py}}]
def plot_convergence(filename):
    data  = np.loadtxt(filename)
    N_vec = data[:, 0].astype(int)
    omega = data[:, 1]          # primera banda

    ref   = omega[-1]
    error = np.abs(omega[:-1] - ref)
    N_sub = N_vec[:-1]

    plt.figure(figsize=(8, 6))
    plt.loglog(N_sub, error, 'o-', label='Error numerico')
    # Referencia de pendiente N^{-2}
    ref_slope = error[0] * (N_sub[0] / N_sub) ** 2
    plt.loglog(N_sub, ref_slope, '--', label=r'$N^{-2}$')
    plt.xlabel('Resolucion $N$')
    plt.ylabel(r'$|\omega_N - \omega_\mathrm{ref}|$')
    plt.legend()
    plt.grid(True, which='both', alpha=0.5)
    out_img = filename.replace('.dat', '.png')
    plt.savefig(out_img, dpi=300, bbox_inches='tight')
    print(f"Imagen guardada: {out_img}")
\end{lstlisting}

% =============================================================================
\section{Verificación y Diagnósticos}
% =============================================================================

\subsection{Residuos ARPACK}

Tras la convergencia, cada autovector verificado satisface:
\begin{equation}
  r_n = \frac{\norm{\Hhat\,\hat{y}_n - \omega_n\,\hat{y}_n}_2}{|\omega_n|}
        < \texttt{tol} = 10^{-10}.
\end{equation}
Los residuos obtenidos en las pruebas ejecutadas son del orden de
$6\times10^{-12}$ a $1.5\times10^{-10}$ para $N=16$, confirmando la
correcta convergencia del solver.

\subsection{Verificación de hermiticidad}

Activando \texttt{-{}-verbosity 2}, el código calcula y reporta:
\begin{equation}
  \delta_H = \max_{i,j} \abs{\hat{H}_{ij} - \conj{\hat{H}_{ji}}}.
\end{equation}
Un valor bien ensamblado produce $\delta_H \lesssim N_\text{dof}\cdot\epsilon_\text{mach}$,
donde $\epsilon_\text{mach} \approx 2.2\times10^{-16}$ es la precisión de máquina.

\subsection{Ortonormalidad modal}

Con el flag \texttt{-{}-verify-orthogonality}, se computa la matriz de Gram:
\begin{equation}
  G_{mn} = \langle \hat{y}_m | \hat{y}_n \rangle = \delta_{mn}
  \quad \Leftrightarrow \quad
  \langle x_m | A | x_n \rangle = \delta_{mn} \quad \text{(ec.\ 13)}.
\end{equation}
Se reportan $\max_{m\neq n}|G_{mn}|$ (debería ser $< 10^{-6}$) y
$\max_m ||G_{mm}| - 1|$ (debería ser $< 10^{-6}$).

% =============================================================================
\section{Conclusiones}
% =============================================================================

Se ha implementado y documentado el solucionador de estructura de bandas
fotónicas \texttt{hermetic-bands}, basado en la formulación Hermitiana
de variable auxiliar de Raman y Fan~\cite{Raman2010}.
Los resultados clave de la auditoría son:

\begin{enumerate}
  \item \textbf{Corrección de dos bugs críticos}: la doble transformación
    espectral en \texttt{zneupd} modo 3 y el uso de \texttt{"LM"} en
    lugar de \texttt{"LR"} en ARPACK producían autovalores incorrectos
    y podían seleccionar modos no físicos del espacio nulo.
    Ambos han sido corregidos.

  \item \textbf{Normalización de unidades}: las constantes físicas
    $\mu_0$, $\varepsilon_0$ y $c$ fueron corregidas a sus valores
    normalizados ($=1$), fundamentales para la correcta construcción
    de la matriz de masa $A$.

  \item \textbf{Hermiticidad por construcción}: $C_H = C_E^\dagger$
    se garantiza mediante \texttt{conj\_transpose}, y las fases de
    Bloch alternan de signo correctamente al transponerse, evitando
    autovalores espurios.

  \item \textbf{LU denso justificado}: para $N \leq 30$, la
    factorización LU densa de LAPACK es la elección óptima: determinista,
    de precisión de máquina y sin dependencias externas adicionales.

  \item \textbf{Separación correcta de pérdidas}: $\Gamma$ entra
    únicamente en $\Hhat_\text{nh}$; el Hamiltoniano lossless $\Hhat_0$
    permanece inalterado.
\end{enumerate}

El proyecto compila y ejecuta correctamente bajo Arch Linux con BLAS,
LAPACK, ARPACK y GFortran del sistema, produciendo residuos ARPACK
del orden de $10^{-12}$--$10^{-10}$ para resoluciones $N=16$--$30$.

% =============================================================================
\begin{thebibliography}{9}
% =============================================================================

\bibitem{Raman2010}
A.~P.~Raman y S.~Fan,
\textit{Photonic Band Structure of Dispersive Metamaterials Formulated as
a Hermitian Eigenvalue Problem},
Phys.\ Rev.\ Lett.\ \textbf{104}, 087401 (2010).
\href{https://doi.org/10.1103/PhysRevLett.104.087401}{DOI:10.1103/PhysRevLett.104.087401}

\bibitem{Taflove2005}
A.~Taflove y S.~C.~Hagness,
\textit{Computational Electrodynamics: The Finite-Difference Time-Domain Method},
3.ª ed., Artech House, Boston (2005).

\bibitem{Lehoucq1998}
R.~B.~Lehoucq, D.~C.~Sorensen y C.~Yang,
\textit{ARPACK Users' Guide: Solution of Large-Scale Eigenvalue Problems
with Implicitly Restarted Arnoldi Methods},
SIAM, Filadelfia (1998).

\bibitem{Anderson1999}
E.~Anderson et al.,
\textit{LAPACK Users' Guide},
3.ª ed., SIAM, Filadelfia (1999).

\bibitem{Joannopoulos2008}
J.~D.~Joannopoulos, S.~G.~Johnson, J.~N.~Winn y R.~D.~Meade,
\textit{Photonic Crystals: Molding the Flow of Light},
2.ª ed., Princeton University Press (2008).

\end{thebibliography}

\end{document}
